\subsection{Существование верхнего и нижнего пределов у любой последовательности.}

\subsubsection*{Определения}
Пусть $P$ — множество всех частичных пределов последовательности $\{a_n\}$ (включая $\pm \infty$).
\begin{itemize}
    \item $L := \overline{\lim_{n\to\infty}} a_n$ — \textbf{верхний предел} (наибольший частичный предел).
    \item $l := \underline{\lim_{n\to\infty}} a_n$ — \textbf{нижний предел} (наименьший частичный предел).
\end{itemize}

\textbf{Теорема.} У любой последовательности $\{a_n\}$ существуют верхний и нижний пределы (конечные или бесконечные).
\[ \forall \{a_n\} \; \exists \overline{\lim_{n\to\infty}} a_n \quad \text{и} \quad \exists \underline{\lim_{n\to\infty}} a_n \]


\begin{tcolorbox}[title=Доказательство (для верхнего предела), breakable]
	Рассмотрим два случая:
	\begin{enumerate}
		\item \textbf{$\{a_n\}$ не ограничена сверху.}
			Это означает, что $\forall M$ найдется бесконечно много членов, больших $M$. Построим подпоследовательность:
			\begin{itemize}
			    \item $a_{n_1} > 1$
			    \item $a_{n_2} > 2, \;\; n_2 > n_1$
			    \item $\dots$
			    \item $a_{n_k} > k, \;\; n_k > n_{k - 1}$
			\end{itemize}
			Тогда $a_{n_k} \to +\infty$. Значит, $+\infty$ — частичный предел.
			Больше чем $+\infty$ быть не может, следовательно:
			\[ \overline{\lim_{n\to\infty}} a_n = +\infty \]

		\item \textbf{$\{a_n\}$ ограничена сверху.}
			Разобьем на два подслучая:
			\begin{enumerate}
				\item \textbf{У $\{a_n\}$ есть конечные частичные пределы.}
					Обозначим $\mathcal{M}$ — множество всех конечных частичных пределов.
					$\mathcal{M} \neq \varnothing$ (по теореме Больцано-Вейерштрасса, если последовательность ограничена и снизу; если нет — см. ниже) и $\mathcal{M}$ ограничено сверху (так как сама $\{a_n\}$ ограничена сверху).
					
					По принципу полноты $\mathbb{R}$ (Вейерштрасса) существует супремум:
					\[ M = \sup \mathcal{M} \]
					Докажем, что $M \in \mathcal{M}$ (т.е. $M$ — сам является частичным пределом).
					В любой $\varepsilon$-окрестности супремума $M$ найдется элемент из $\mathcal{M}$ (какой-то частичный предел $a'$). В окрестности $a'$ лежит бесконечно много членов $a_n$. Значит, и в окрестности $M$ лежит бесконечно много членов $a_n$.
					Следовательно, $M$ — частичный предел.
					Так как $M$ — супремум, то это максимальный частичный предел:
					\[ M = \overline{\lim_{n\to\infty}} a_n \]

				\item \textbf{У $\{a_n\}$ нет конечных частичных пределов.}
					В этом случае докажем, что $a_n \to -\infty$.
					Предположим противное: $a_n \not\to -\infty$.
					Тогда существует $E > 0$, такое что для любого $N$ найдется $n > N$, где $a_n \ge -E$.
					
					Таким образом, можно выделить подпоследовательность, ограниченную снизу числом $-E$. Так как мы в случае 2, она ограничена и сверху.
					По теореме Больцано-Вейерштрасса из ограниченной подпоследовательности можно выделить сходящуюся.
					Значит, существует конечный частичный предел. \textbf{Противоречие} с условием подпункта (b).
					
					Следовательно, $a_n \to -\infty$, и тогда:
					\[ \overline{\lim_{n\to\infty}} a_n = -\infty \]
			\end{enumerate}
	\end{enumerate}
\end{tcolorbox}
